\documentclass[11pt]{article}
\usepackage[margin=1in]{geometry}
\usepackage{amsmath,amssymb,amsthm}
\usepackage{graphicx}
\usepackage[hidelinks]{hyperref}

\newtheorem{theorem}{Theorem}
\newtheorem{remark}{Remark}

\newcommand{\GammaSet}{\Gamma}
\newcommand{\pt}{\mathrm{pt}}
\newcommand{\abs}[1]{\left\lvert #1 \right\rvert}

\graphicspath{{runs/fa_banach_001/solutions/full/2411.09430_pointwise_least_locally_solid_topology/}}

\title{Pointwise Convergence Is the Least Locally Solid Hausdorff Topology}
\author{FA Banach run \texttt{fa\_banach\_001}, agent \texttt{agent\_lane\_06}}
\date{2026-06-19}

\begin{document}
\maketitle

\section*{Status}

This packet records a candidate full solution, likely valid, to Question 6.c.1
of Kiwerski and Kolwicz \cite{KK}: whether local convergence in measure is the
coarsest locally solid Hausdorff topology in the class of Banach sequence
spaces.

\section*{Source Passage}

\begin{center}
\includegraphics[width=0.98\linewidth]{figures/open_problem_crop.png}
\end{center}

\section*{Transcription of the Open Question}

The source paper asks:
\begin{quote}
``Is the topology of local convergence in measure the coarsest locally solid
Hausdorff topology in the class of Banach sequence spaces?''
\end{quote}

Immediately before the question, the authors note that for a Banach sequence
space the topology of pointwise convergence coincides with local convergence in
measure.  They also note Labuda's 1987 theorem for Orlicz spaces
\cite{Labuda}.

\section*{Definitions Used}

Let \(\GammaSet\) be countable and equipped with counting measure.  Following
\cite{KK}, a Banach sequence space \(X\) over \(\GammaSet\) is a Banach space
which is a linear subspace of \(\omega(\GammaSet)\), contains
\({\bf 1}_F\) for every finite \(F\subseteq\GammaSet\), and has the ideal
property: if \(\abs{x(\gamma)}\le \abs{y(\gamma)}\) for all
\(\gamma\in\GammaSet\) and \(y\in X\), then \(x\in X\).

For \(\gamma\in\GammaSet\), write \(e_\gamma={\bf 1}_{\{\gamma\}}\).  Let
\(\tau_{\pt}\) denote the topology of pointwise convergence on \(X\), that is,
the initial topology generated by the coordinate maps
\[
  \pi_\gamma:X\to\mathbb R,\qquad \pi_\gamma(x)=x(\gamma).
\]

A subset \(U\subseteq X\) is solid if \(x\in U\), \(y\in X\), and
\(\abs{y(\gamma)}\le \abs{x(\gamma)}\) for all \(\gamma\) imply \(y\in U\).
A vector topology on \(X\) is locally solid if zero has a base of solid
neighborhoods.

\section*{Theorem}

\begin{theorem}
Let \(X\) be a Banach sequence space over a countable set \(\GammaSet\).  Then
the pointwise topology \(\tau_{\pt}\) is the least Hausdorff locally solid
vector topology on \(X\).  Consequently, for the counting measure on
\(\GammaSet\), the topology of local convergence in measure is the coarsest
locally solid Hausdorff topology in the class of Banach sequence spaces.
\end{theorem}

\section*{Proof Intuition}

The only possible way for a Hausdorff vector topology to be too coarse is for
it not to see some coordinate.  Solidity prevents this.  If a solid
neighborhood contains \(x\), then it also contains the one-coordinate fragment
\(x(\gamma)e_\gamma\).  Hausdorffness on the one-dimensional coordinate axis
forces these fragments to have small scalar coefficient inside sufficiently
small neighborhoods.  Hence every coordinate map is continuous for every
Hausdorff locally solid topology.

\section*{Proof}

First observe that \(\tau_{\pt}\) is itself a Hausdorff locally solid vector
topology.  Indeed, a base at zero is given by the finite-coordinate sets
\[
  U(F,\varepsilon)
  =
  \{x\in X:\abs{x(\gamma)}<\varepsilon
       \text{ for all }\gamma\in F\},
\]
where \(F\subseteq\GammaSet\) is finite and \(\varepsilon>0\).  These sets are
solid, and the resulting topology is Hausdorff because the coordinates
separate points of \(X\subseteq\omega(\GammaSet)\).

Now let \(\tau\) be any Hausdorff locally solid vector topology on \(X\).  We
prove that \(\tau_{\pt}\subseteq\tau\).  Fix \(\gamma\in\GammaSet\) and
\(\varepsilon>0\).  Since \(e_\gamma\ne 0\) and \(\tau\) is Hausdorff, there is
a balanced \(\tau\)-neighborhood \(N\) of zero such that
\(e_\gamma\notin N\).  Put \(U=\varepsilon N\).  If
\(\lambda e_\gamma\in U\), then \((\lambda/\varepsilon)e_\gamma\in N\).  Were
\(\abs{\lambda}\ge\varepsilon\), balancedness of \(N\) would imply
\[
  e_\gamma
  =
  \frac{\varepsilon}{\lambda}
  \left(\frac{\lambda}{\varepsilon}e_\gamma\right)
  \in N,
\]
a contradiction.  Hence
\[
  U\cap \mathbb R e_\gamma
  \subseteq
  \{\lambda e_\gamma:\abs{\lambda}<\varepsilon\}.
\]

Because \(\tau\) is locally solid, choose a solid \(\tau\)-neighborhood
\(V\subseteq U\).  If \(x\in V\), then \(x(\gamma)e_\gamma\in X\) and
\[
  \abs{x(\gamma)e_\gamma(\eta)}
  \le \abs{x(\eta)}
  \qquad(\eta\in\GammaSet).
\]
By solidity, \(x(\gamma)e_\gamma\in V\subseteq U\), so the previous paragraph
gives \(\abs{x(\gamma)}<\varepsilon\).  Thus
\[
  V\subseteq \{x\in X:\abs{\pi_\gamma(x)}<\varepsilon\},
\]
which proves that \(\pi_\gamma\) is \(\tau\)-continuous at zero, hence
continuous.

All coordinate maps are therefore \(\tau\)-continuous.  Since
\(\tau_{\pt}\) is the initial topology generated by these coordinate maps, it
follows that \(\tau_{\pt}\subseteq\tau\).  As \(\tau\) was arbitrary,
\(\tau_{\pt}\) is the least Hausdorff locally solid vector topology on \(X\).

Finally, for the counting measure on a countable set, local convergence in
measure is convergence in measure on finite subsets of \(\GammaSet\).  Since
finite subsets carry counting measure with finitely many atoms, this is
equivalent to coordinatewise convergence on every finite set, i.e. to
\(\tau_{\pt}\).  The stated affirmative answer follows.
\(\square\)

\section*{Verification Notes}

No numerical verification is involved.  The proof uses only the structural
axioms in the source paper's definition of Banach sequence spaces: finite
coordinate vectors belong to \(X\), and the ideal property guarantees that
coordinate fragments \(x(\gamma)e_\gamma\) belong to \(X\) and lie in solid
neighborhoods whenever \(x\) does.

The proof does not use the norm topology.  It gives the stronger statement
that pointwise convergence is below every Hausdorff locally solid vector
topology on the underlying sequence lattice.

\section*{Novelty and Literature Check}

The run indexes were searched for arXiv:2411.09430, ``Kadets'', ``Kadec'',
``Klee'', ``direct sums'', ``coarsest locally solid'', ``minimal locally
solid'', ``local convergence in measure'', and related phrases.  No existing
packet or attempt addresses this question.

External searches on 2026-06-19 for exact phrases from Question 6.c.1 and for
``coarsest locally solid Hausdorff topology'' with ``Banach sequence space''
found the source paper and adjacent minimal-topology literature, especially
Kandic and Taylor \cite{KT}, but no later work explicitly resolving this
Banach sequence space question.  Labuda's theorem \cite{Labuda}, cited by the
source paper, treats Orlicz spaces and is consistent with the argument here.

Human review should focus on the standard locally solid topology convention:
zero has a base of solid neighborhoods.  Under that convention, the
coordinate-axis argument above is the complete proof.

\begin{thebibliography}{9}

\bibitem{KK}
T. Kiwerski and P. Kolwicz,
\emph{Direct sums and abstract Kadets--Klee properties},
arXiv:2411.09430.

\bibitem{Labuda}
I. Labuda,
\emph{Submeasures and locally solid topologies on Riesz spaces},
Mathematische Zeitschrift 195 (1987), 179--196.

\bibitem{KT}
M. Kandic and M. A. Taylor,
\emph{Metrizability of minimal and unbounded topologies},
arXiv:1709.05407.

\end{thebibliography}

\end{document}
